\chapter{Estilos arquitectónicos y sistemas distribuidos}

\begin{objetivos}
El lector podrá comparar monolito modular, microservicios, arquitectura orientada a eventos y serverless; identificar costos de distribución; y diseñar una evolución incremental basada en señales observables.
\end{objetivos}

\section{Los estilos son restricciones deliberadas}
Un estilo arquitectónico define un conjunto de elementos, relaciones y restricciones. Capas controla la dirección de dependencias; hexagonal separa política de adaptadores; microservicios agrega autonomía de despliegue y datos; eventos desacopla temporalmente productores y consumidores. Combinar estilos es normal, pero cada combinación añade reglas que deben verificarse.

\section{Monolito modular}
Un monolito modular se despliega como una unidad y contiene límites internos explícitos. Sus ventajas son depuración, transacciones y operación simples. Sus riesgos aparecen cuando los límites se ignoran, la base de datos se convierte en interfaz informal y cada cambio exige comprender todo el sistema.

Para un equipo pequeño y un dominio en aprendizaje, suele ser un punto de partida prudente. La modularidad crea opciones: un módulo puede extraerse después si existe evidencia de que necesita escalar, desplegarse o gobernarse de forma independiente.

\section{Microservicios}
Un microservicio posee una responsabilidad coherente, despliegue independiente y control de sus datos. Dividir por entidades produce servicios anémicos y conversaciones excesivas. El límite debe reflejar capacidad de negocio, equipo y ciclo de cambio.

La distribución introduce latencia, fallos parciales, versionamiento de contratos, observabilidad, seguridad entre servicios, consistencia eventual y automatización. Si dos servicios siempre cambian y se despliegan juntos, la separación puede ser accidental.

\section{Eventos y mensajería}
Un evento declara un hecho ocurrido: \software{ReservaConfirmada}. No debe ser una orden disfrazada ni incluir el modelo interno completo. Los consumidores pueden reaccionar de forma independiente, pero se deben manejar duplicados, orden, reintentos y evolución del esquema.

\section{Serverless}
Serverless delega infraestructura a una plataforma y factura por uso o capacidad. Puede reducir operación para cargas intermitentes, pero introduce límites de ejecución, arranques, observabilidad distribuida y dependencia del proveedor. Las alternativas libres, como OpenFaaS o Knative, también requieren operar una plataforma subyacente; por tanto, no eliminan el costo, lo redistribuyen.

\section{Árbol de decisión}
\begin{figure}[H]
\centering
\begin{tikzpicture}[scale=0.88,transform shape,node distance=1.05cm,every node/.style={align=center,font=\small}]
  \node[draw=AzulUD,fill=AzulClaro,rounded corners] (start) {¿Existe necesidad demostrable\\de despliegue independiente?};
  \node[draw=Verde,fill=VerdeClaro,rounded corners,below left=of start] (mono) {No: fortalecer\\monolito modular};
  \node[draw=AzulUD,fill=AzulClaro,rounded corners,below right=of start] (team) {Sí: ¿hay equipo y\\plataforma operativa?};
  \node[draw=Naranja,fill=NaranjaClaro,rounded corners,below left=of team] (prep) {No: crear capacidades\\antes de distribuir};
  \node[draw=Verde,fill=VerdeClaro,rounded corners,below right=of team] (svc) {Sí: extraer una\\capacidad delimitada};
  \draw[-{Latex},thick] (start) -- node[above left]{no} (mono);
  \draw[-{Latex},thick] (start) -- node[above right]{sí} (team);
  \draw[-{Latex},thick] (team) -- node[above left]{no} (prep);
  \draw[-{Latex},thick] (team) -- node[above right]{sí} (svc);
\end{tikzpicture}
\caption{Decisión conservadora antes de distribuir.}
\label{fig:decision-distribucion}
\end{figure}

\section{Patrones de evolución}
\begin{itemize}
  \item \textbf{Strangler}: enrutar gradualmente capacidades hacia un componente nuevo.
  \item \textbf{Anti-corruption layer}: traducir modelos para proteger un dominio de un legado.
  \item \textbf{Outbox}: guardar cambio de negocio y mensaje en una transacción local, publicándolo después.
  \item \textbf{Saga}: coordinar transacciones locales con compensaciones; no restaura automáticamente el estado exacto anterior.
  \item \textbf{Circuit breaker}: detener llamadas a una dependencia fallida para permitir recuperación; no sustituye límites de tiempo.
\end{itemize}

\begin{ejercicio}[title={Laboratorio. Simulación de fallos}]
Modele la confirmación de una reserva y el envío de correo como dos procesos. Simule que el correo falla en 30\% de los intentos. Compare: transacción síncrona, reintento inmediato, cola con reintento y outbox. Registre reservas confirmadas, correos duplicados y tiempo total.
\end{ejercicio}

\begin{validacion}
La estrategia propuesta debe demostrar: ninguna reserva confirmada se pierde; un fallo de correo no revierte indebidamente la reserva; los mensajes duplicados no producen efectos duplicados; los fallos agotados quedan visibles para operación.
\end{validacion}

\begin{ejercicio}[title={Ejercicio de arquitectura}]
Seleccione un módulo candidato a extracción. Defina señal de activación, contrato, propiedad de datos, plan de migración, observabilidad y estrategia de reversión. Si no puede identificar una señal medible, argumente por qué debe permanecer en el monolito.
\end{ejercicio}

\begin{riesgo}
Microservicios no son una medida de madurez. La madurez consiste en poder justificar límites, operar fallos y cambiar con seguridad. Distribuir antes de dominar pruebas, entrega y observabilidad amplifica los problemas existentes.
\end{riesgo}
